\clearpage
\renewcommand{\sectioncolor}{quantum}
\section{After sequencing: the look-alike gene trap}
\markboth{The paralogue trap}{}

\noindent\textbf{\color{quantum}The biology question:} \emph{when BLAST returns an 82\,\% hit to
lukF-PV, is the gene there or not?}

\noindent This is the one place in our pipeline where a mathematical mistake has \emph{already}
produced a wrong biological conclusion --- documented in our own \texttt{METHODS\_v1.7.md} §4.6 as
``the incorrect PVL inference in report v1.0''.

\subsection{What the data actually looks like}

\begin{center}
\begin{tikzpicture}[font=\footnotesize]
% axis: identity 38..102 mapped to x 0..14.2
\def\X(#1){(#1-38)/64*14.2}
\draw[slate!55,line width=1.2pt,-{Stealth[length=5pt]}] (0,0)--(14.7,0);
\foreach \i in {40,50,60,70,80,90,100}{
  \pgfmathsetmacro{\xx}{(\i-38)/64*14.2}
  \draw[slate!55] (\xx,-0.12)--(\xx,0.12);
  \node[font=\scriptsize,below=3pt] at (\xx,-0.1) {\i};}
\node[font=\scriptsize,text=slate,anchor=north] at (7.1,-1.85)
   {BLAST nucleotide identity to the reference gene (\%)};
% threshold lines, labels BELOW
\foreach \t/\lab/\c in {80/{80\,\% floor\\ResFinder}/quantum, 90/{90\,\% floor\\VFDB}/quantum}{
  \pgfmathsetmacro{\xx}{(\t-38)/64*14.2}
  \draw[\c,line width=1pt,dash pattern=on 3pt off 2pt] (\xx,-0.72)--(\xx,0.95);
  \node[font=\scriptsize,text=\c,anchor=north,align=center] at (\xx,-0.78) {\lab};}
% real calls
\foreach \i in {100,100,100,99.95,99.88,99.85,99.59,99.48,99.10,98.42}{
  \pgfmathsetmacro{\xx}{(\i-38)/64*14.2}
  \fill[bio] (\xx,0) circle (2.4pt);}
\pgfmathsetmacro{\xg}{(99.2-38)/64*14.2}
\node[font=\scriptsize\bfseries,text=bio,anchor=south east,align=right] at (\xg+0.55,1.28)
  {\textbf{every reported call}\\98.42 -- 100 \%};
\draw[bio,line width=0.7pt] (\xg,0.20)--(\xg,1.22);
% lukF-PV
\pgfmathsetmacro{\xl}{(82.31-38)/64*14.2}
\fill[alert] (\xl,0) circle (3.0pt);
\node[font=\scriptsize\bfseries,text=alert,anchor=south,align=center] at (\xl,0.38)
  {\textbf{lukF-PV} 82.31 \%\\{\scriptsize 98.8 \% of the gene covered}};
% gyrA/parC
\pgfmathsetmacro{\xc}{(41-38)/64*14.2}
\fill[ember] (\xc,0) circle (3.0pt);
\node[font=\scriptsize\bfseries,text=ember,anchor=south west,align=left] at (\xc-0.05,0.38)
  {gyrA$\leftrightarrow$parC\\{\scriptsize cross-hit $\approx$41 \%}};
\end{tikzpicture}
\end{center}
\vspace{4pt}
\begin{center}\begin{minipage}{0.93\textwidth}\footnotesize\color{slate}
\textbf{Figure 7.}\; The identity landscape from \texttt{METHODS\_v1.7.md} §5. The 90\,\% VFDB
threshold was chosen \emph{specifically} to keep the leukocidin paralogue out. It works --- but it
works because somebody already knew the answer.
\end{minipage}\end{center}

\begin{alertbox}[title={Why the threshold is the wrong tool, even though it gave the right answer}]
Notice what saved us. It was \textbf{not} the 82.31\,\% number. It was that the same region hits
\texttt{hlgB} --- gamma-haemolysin, the known $\approx$80\,\%-identity paralogue of \texttt{lukF-PV}
--- at \textbf{99.90\,\%}. The reasoning was:
\begin{center}\itshape
this looks like lukF-PV, but it looks far more like hlgB, and hlgB is definitely present.
\end{center}
That is a \textbf{comparison between two hypotheses}, not a comparison against a fixed number. And
comparison between two hypotheses has a standard mathematical form:
\[
\Lambda \;=\;
\frac{P(\text{this alignment}\mid \text{the gene is really }\texttt{lukF-PV})}
     {P(\text{this alignment}\mid \text{it is really the paralogue }\texttt{hlgB})}
\]
We computed $\Lambda$ by hand, in prose, once. \textbf{It should be computed by the script, for every
call, automatically.}
\end{alertbox}

\begin{plain}
Percent identity throws away the most informative thing in the alignment: \emph{which} positions
differ. Two genes can both be 95\,\% identical to a reference and yet one of them differs only at
positions that every member of the family varies at, while the other differs at the conserved
catalytic residues. Percent identity scores them the same. A profile --- a model of which positions
are allowed to vary --- does not.
\end{plain}

\subsection{Three replacements, in increasing order of effort}

\begin{center}\small
\begin{tabular}{@{}>{\raggedright\arraybackslash}p{3.7cm}>{\raggedright\arraybackslash}p{6.2cm}>{\raggedright\arraybackslash}p{6.2cm}@{}}
\toprule
\rowcolor{quantum!13}
\textbf{\color{quantum}Method} & \textbf{\color{quantum}What it does} & \textbf{\color{quantum}Effort / payoff}\\
\midrule
\textbf{Best-vs-second-best ratio} \newline (do this first) &
 For every call, also report the best hit to a \emph{different} gene family, and the ratio. Our
 \texttt{06\_pvl\_test.py} already does this for PVL --- generalise it to all five databases. &
 \textbf{A few hours.} Removes the entire class of error that produced the v1.0 mistake. Requires no
 new software.\\
\textbf{Profile HMMs} \newline (HMMER) &
 Score against a \emph{model of the family} rather than one sequence. Positions that are conserved
 across the family are weighted more heavily. &
 A day, plus building profiles. This is what AMRFinderPlus already uses internally.\\
\textbf{Protein language model embeddings} \newline (ESM-2) &
 Represent each protein as a learned vector; orthologues cluster, paralogues separate --- even at
 identical percent identity. &
 Genuinely better for hard families. See \S8 for the caution that comes with it.\\
\bottomrule
\end{tabular}
\end{center}

\begin{dobox}[title={Why this matters beyond one gene}]
Our Limitations section currently has to defend a threshold. With a likelihood ratio it would not
have to: the number would carry its own justification. And the same machinery covers the
\texttt{gyrA}/\texttt{parC} cross-hit at $\approx$41\,\%, which we currently exclude with a hand-set
$\ge$90\,\% orthologue filter --- another threshold that happens to be right.
\end{dobox}
